\begin{table}[htbp]
\centering
\footnotesize
\caption{Moderate 36-month rolling policy versus a risk-compatible benchmark.}
\label{tab:gbi-moderate-rolling-benchmark}
\begin{tabular}{lrrr}
\toprule
Metric & Rolling policy & Benchmark & Difference \\
\midrule
Annualized return & 9.11% & 9.45% & -0.34% \\
Annualized volatility & 7.38% & 9.36% & -1.97% \\
Sharpe ratio & 1.22 & 1.02 & +0.21 \\
Maximum drawdown & -6.85% & -14.43% & +7.58% \\
Historical ES (95\%) & 3.61% & 4.89% & -1.28% \\
\bottomrule
\end{tabular}
\vspace{2pt}
\begin{minipage}{0.96\linewidth}
\footnotesize
Notes: Monthly out-of-sample returns after a 36-month estimation period. Both portfolios rebalance quarterly with 10bp one-way transaction costs. Sharpe assumes a zero risk-free rate. For return and Sharpe, a positive difference favors the rolling policy; for volatility and ES, a negative difference favors it; for maximum drawdown, a positive difference means a smaller loss.
\end{minipage}
\end{table}
